%=====================================================================
%  ch06 —— 原文 1794–2231：§3.4 Hidden Symmetries
%=====================================================================

\subsection{隐藏的对称性}

\ssszh 上一小节我们看到，局部化满足一条重要的相容性——它把等变 D-模与
Harish-Chandra 模互换。把这一断言展开之后，它归结为关于函子形式的一个相当具体而
初等的观察。

然而必须认识到，这样的陈述，尤其是在它刚被提出的时候，并不容易嵌入表示论中已知的
一般形式体系。

现在我们要说明，这条相容性其实是局部化两侧都存在的一个更丰富结构的影子。在几何
一侧，进一步的对称性即使不算显然，也并不完全出人意料。然而在表示论一侧，它就远不那么
显然了。可以恰当地把它看作李代数表示范畴的一个隐藏对称性。不过正如我们将要说明的，
它的各个片段曾在个案基础上被看到并被卓有成效地使用，但完整的结构在几十年间都没有浮现。

\ssszh 我们的出发点将是下面的类比：
%
\[
   \Dmod(\sB) \quad \text{与} \quad \Fun(\sB(\Fq)).
\]
%
回忆后者带有 $G(\Fq)$ 的作用，这是 $G(\Fq)$ 在 $\sB(\Fq)$ 本身上作用的线性化。

这提示了如下的想法——我们应当把 $\Dmod(\sB)$ 看成群 $G$ 的某种表示，同样是把 $G$ 在
$\sB$ 上的作用线性化。关键在于，局部化等价
%
\[
    \Dmod(\sB) \simeq \fg\mod_0
\]
%
应当是 $G$-模的等价，特别地，李代数表示的范畴本身应当是一个 $G$-模。要把这一点说
精确，需要对代数群在范畴（而非向量空间）上的作用有一个好的概念。

\ssszh 我们先在第 \ref{s:kandx} 小节那样的较一般设定（$K$ 与 $X$）中，对这个想法作
一个合理性检查。

回忆局部化与整体截面函子
%
\[
     \on{Loc}: \fk\mod \rightleftarrows \Dmod(X): \Gamma(X, -).
\]
%
无论 $K$ 在范畴 $\sC$ 上的作用该是什么，至少对 $K$ 的每个元素 $g$，我们应当有范畴的
自同构
%
$$
    g: \sC \xrightarrow{\sim}\sC,
$$
%
而且它们在复合下应当适当相容。现在我们来为 $\fk\mod$ 与 $\Dmod(X)$ 找出这样的
自同构，使局部化把两者交织起来。

我们从几何一侧开始。由于 $X$ 是 $K$-簇，用 $g$ 作用给出簇本身的自同构
$$g: X \xrightarrow{\sim} X$$，从而给出其 D-模范畴的自同构
%
\[
      g_*: \Dmod(X) \xrightarrow{\sim} \Dmod(X).
\]
%
例如，对 $X$ 的任意点 $x$，它把 $\delta_x$ 送到 $\delta_{gx}$。注意对 $K$ 中任意
$g$ 与 $g'$，有典范同构
%
\[
     g_* \circ g'_* \simeq (gg')_*,
\]
%
即群的乘法与自函子的复合之间具有相容性。

为看出 $\fk\mod$ 上相应的算子应当是什么，我们来看它对整体截面的作用。注意对 $X$ 上
任一 D-模 $M$，
%
\[
   \Gamma(X, M) \quad \text{与} \quad \Gamma(X, g_* M)
\]
%
的底向量空间有典范的等同，但整体微分算子的作用被 $g$ 扭了一下。

有鉴于此，记
%
\[
       \on{Ad}_{g, *}: \fk\mod \xrightarrow{\sim} \fk\mod
\]
%
为下述自同构：它把作用为 $\fk \otimes N \rightarrow N$、$X \otimes n \mapsto X
\cdot n$ 的 $\fk$-模 $N$，送到底向量空间相同、作用为 $X \otimes n \mapsto
\on{Ad}_{g^{-1}}(X) \cdot n$ 的 $\fk$-模。这里同样有相容性
%
\[
   \on{Ad}_{g, *} \circ \on{Ad}_{g', *} \simeq \on{Ad}_{gg', *}.
\]
%

\ssszh 由上述讨论，对 $K$ 中每个 $g$，我们得到典范地交换的图
%
\[
\xymatrix{\Dmod(X) \ar[d]^{g_*} \ar[rr]^{\Gamma(X, -)} && \fk\mod
\ar[d]^{\on{Ad}_{g, *}} \\ \Dmod(X) \ar[rr]^{\Gamma(X, -)} & &
\fk\mod. }
\]
%
也就是说，$\Gamma(X, -)$ 把目前为止在两个范畴上构造的 $K$ 作用交织起来。读者不妨
验证 $\on{Loc}$ 也有同样的性质。

\ssszh 现在到了重要的一点。我们至今为止期望一个代数群的范畴表示所具有的结构量，
显然还不够。例如，我们只用到它的闭点，并把它们看作抽象群。

为继续推进，回忆 $G(\Fq)$ 在 $\Fun(\sB(\Fq))$ 上的作用可延拓为群代数
$\Fun(G(\Fq))$ 的作用。按照函数-层对应，我们发现导出范畴 $\Dmod(K)$ 在卷积下是
幺半的。具体地说，给定 $\Dmod(K)$ 中两个层 $\sF$ 与 $\sG$，其卷积为
%
\begin{equation} \label{e:convalgdmod}
  \sF \star \sG := m_*\big( p_1^! \sF \overset{!}{\otimes} p_2^! \sG \big),
\end{equation}
%
其中用 $m, p_1, p_2: K \times K \rightarrow K$ 分别记乘法映射与到第一、第二因子的
投影，并对态射使用显然的函子性。

我们强调，在函数-层对应下，公式 \eqref{e:convalgdmod} 与群代数中通常的乘法相符。
与此相关，幺半单位是支撑在群单位元处的 delta 函数 $\delta_e$。有点启发式地说，正如
通常的群代数那样，可以把一般的 D-模想成 delta D-模的「连续线性组合」。

类似地，$\Dmod(X)$ 自然地成为 $\Dmod(K)$ 的模范畴。具体地说，给定
$\Dmod(K)$ 中的层 $\sF$ 与 $\Dmod(X)$ 中的 $\sM$，其卷积为
%
\begin{equation} \label{e:conv}
     \sF \star \sM := a_*\big( p_1^! \sF \overset{!}{\otimes} p_2^! \sM\big),
\end{equation}
%
其中 $a, p_1, p_2: K \times X \rightarrow X$ 记作用映射与投影映射。

按构造，与支撑在 $K$ 的点 $g$ 处的 delta 函数作卷积，重现已讨论过的对称性，即有
函子的典范同构
%
\begin{equation}
\label{e:action}
 \delta_g \star (-) \simeq g_*: \Dmod(X) \rightarrow \Dmod(X).
\end{equation}
%
非正式地说，我们因此把作用从单个 delta 函数延拓到了它们的全部线性组合，即
$\Dmod(K)$。

下一个任务是为李代数表示找出 $\Dmod(K)$ 的相应作用。在此之前，我们想对 D-模一侧
再作几点评论。

\begin{re} 注意，关于自同构 $g_*$ 与 $\on{Ad}_{g, *}$ 的讨论在阿贝尔范畴或导出范畴
中都同样可行；但要使这个幺半结构按所希望的方式运作，即甚至要有一个像样的 D-模前推，
就必须在导出范畴中工作。目前三角范畴已经够用，但下一节我们将被迫要求更多。

还要注意，要在卷积下封闭，我们必须立刻离开有限生成 D-模的世界。原型例子是 $\sD_K$
与自身的卷积，可以证明
%
\[
    \sD_K \star \sD_K \simeq \sD_K \underset{\underline{k}}{\otimes} \OO_K.
\]
%
\end{re}

\begin{re} 把卷积与作用的公式改写为
%
\begin{align*}
&\Dmod(K) \otimes \Dmod(K) \simeq \Dmod(K \times K) \xrightarrow{m_*}
\Dmod(K) \quad \text{与} \\ &\Dmod(K) \otimes \Dmod(X) \simeq
\Dmod(K \times X) \xrightarrow{a_*} \Dmod(X),
\end{align*}
%
会更有启发，正如函数情形那样，参看公式 \eqref{e:convfun} 与 \eqref{e:actionfun}。
不过为此需要导出范畴的张量积，而这只能在 dg 范畴的层次上进行，见第
\ref{s:catrepgrp} 节。
\end{re}

\ssszh\label{ss:weakinv} 最后，转向 $K$ 在 $\fk\mod$ 上的范畴作用。首先注意，
目前这个作用的构造不像 $\Dmod(X)$ 的情形那样直接。确实，那里我们只需要 D-模作用在
D-模上，这可以用几何函子性；这里我们却需要 D-模作用在李代数表示上，两者表面上相距
更远。

不过，我们可以把范畴 $\fk\mod$ 改写为 D-模的范畴来推进。回忆若 $K$ 作用在簇 $X$
上，则微分算子层 $\sD_X$ 不是等变 D-模。然而它是\emph{弱等变} D-模的原型例子——
这种对象的阿贝尔范畴记作
%
\[
     \Dmod(X)^{K, w, \heartsuit}
\]
%
其定义与第 \ref{s:rigimon} 小节完全相同，只是不加关键的条件 (ii)。

若取 $X = K$，看作右乘下的 $K$-簇，我们断言有典范等价
%
\begin{equation} \label{e:weakinvdmod}
    \Dmod(K)^{K, w, \heartsuit} \simeq \fk\mod^\heartsuit.
\end{equation}
%
为方便读者，我们勾勒证明。首先验证左端有投射生成元 $\sD_K$。这来自它余表示
$K$-不变的整体截面，即 $\Gamma(K, -)^K$，再结合 $K$ 上 $K$-等变拟凝聚层与
$\on{Vect}$ 之间的重言同构。其次，回忆右不变微分算子与包络代数的等同，即
%
\[
     \Gamma(K, \sD_K)^K \simeq U(\fk).
\]
%
总之，我们找到了一个自同态为 $U(\fk)$ 的投射生成元，由 Morita 理论即得结论。

公式 \eqref{e:weakinvdmod} 的等价可延拓为导出范畴的等价
%
\[
    \Dmod(K)^{K, w} \simeq \fk\mod,
\]
%
而前者通过左卷积带有 $\Dmod(K)$ 的作用，所需作用的构造至此完成。

\ssszh 有了这些，局部化与整体截面函子所希望的升级如下。

\begin{theo}[Beilinson--Drinfeld, \cite{bdh}]\label{t:loc+equiv} 设 $X$ 与 $K$ 如
第 \ref{s:kandx} 小节。有 $\Dmod(K)$-等变函子的典范伴随对
%
\[
    \on{Loc}: \fk\mod \rightleftarrows \Dmod(X): \Gamma(X, -).
\]
%
\end{theo}

特别地，定理 \ref{t:bbloc} 中所述的 Beilinson--Bernstein 局部化，是
$\Dmod(G)$ 表示的等价。\footnote{细心的读者会注意到，该局部化定理涉及的不是整个
$\fg\mod$，而是 $\fg\mod_0$。$\Dmod(G)$ 在 $\fg\mod_0$ 上的作用，形式地可由
$Z(\fg)$ 与 $\Dmod(G)$ 在 $\fg\mod$ 上作用的适当交换性推出，也可像第
\ref{ss:weakinv} 小节那样手工构造。}定理 \ref{t:loc+equiv} 的解释，包括对
$\Dmod(K)$-等变性的讨论，我们推迟到第 \ref{s:catrepgrp} 节。

\ssszh 下面指出这一理论的动机性应用。为定向，考虑有限群的子群，不妨暗示性地记作
$H(\Fq) \subset K(\Fq)$。回忆对 $K(\Fq)$ 的任一表示 $\pi$，其
$H(\Fq)$-不变量 $\pi^{H(\Fq)}$ 带有卷积代数
$$\Fun(H(\Fq) \backslash K(\Fq) /H(\Fq))$$ 的作用。

类似地，给定 $K$ 的代数子群 $H$，对 $K$ 的任一范畴表示 $\sC$ 可以取其 $H$-不变
对象 $\sC^H$，它们将带有卷积范畴
%
\[
   \Dmod(H \backslash K  / H)
\]
%
的作用。在目前的例子中，我们得到定理 \ref{t:heqv} 及其导出变体的如下升级。

\begin{theo}[Beilinson--Drinfeld \cite{bdh}]\label{t:lochchmod+eq}存在
$\Dmod(H \backslash K / H)$-等变函子的典范伴随对
%
\[
\on{Loc}: \fk\mod^H \rightleftarrows \Dmod(X)^H: \Gamma(X, -).
\]
%
\end{theo}

一旦形式体系就位，这可以由定理 \ref{t:loc+equiv} 施加 $H$-不变量得到。特别地，
范畴不变量的一般概念（将在第 \ref{s:catrepgrp} 节讨论）在李理论一侧重现
Harish-Chandra 模的范畴
%
\[
   (\fk, H)\mod \simeq \fk\mod^H
\]
%
而在几何一侧重现等变 D-模的范畴。

\ssszh\label{ss:hitchinhecke} 定理 \ref{t:lochchmod+eq} 是 Beilinson--Drinfeld
在他们关于 Hitchin 系统量子化的奠基性工作 \cite{bdh} 中证明的。\footnote{严格说来，
他们需要并证明的是某个无穷维的类似物。}在这个动机问题中，与
$\Dmod(H \backslash K / H)$ 作用的相容性至关重要。即，所希望的结果是 Hecke 特征层的
存在，那是一类关于这个作用具有等变性质的 D-模。它们经由局部化从李代数表示构造
出来，而那些表示被直接证明具有所需性质。

\ssszh 下面是一个不那么精巧的例子。考虑最高权表示的范畴
%
\[
    \fg\mod^B.
\]
%
由上述形式体系，它重言地带有 Hecke 范畴
%
\[
   \Dmod(B \backslash G / B).
\]
%
的作用。我们要强调，在其历史发展中，最高权模与旗簇上 D-模之间的联系是通过局部化
定理才被看到的，特别地，它高度不显然。然而范畴作用的形式体系使它们的关系一目了然。
而且，这迅速而概念性地重现了该范畴的若干重要自函子。

例如，对有限 Weyl 群的每个元素 $w$，可附上某些自函子
%
\[
  j_{w, !} \quad \text{与} \quad j_{w, *}: \fg\mod^B \rightarrow \fg\mod^B,
\]
%
分别称为 Enright 完备化函子与 Arkhipov 扭化函子 \cite{enright79}、\cite{bb82}、
\cite{ark}。几何上，它们由与
%
\[
     BwB \subset G
\]
%
上常值 D-模的 $!$-和 $*$-扩张作卷积给出，其许多基本性质都有简单的几何证明。

作为相关的例子，在 Schubert 簇相交上同调的若干非平凡事实之下，
Kazhdan--Lusztig 猜想可由下述断言推出：
%
\begin{equation} \label{e:brungus}
  \fg\mod_0^B
\end{equation}
%
与 $\Dmod(B \backslash G / B)$ 典范等价。这一点最初是通过建立局部化定理证明的。不过，
利用 Hecke 范畴在 \eqref{e:brungus} 上的作用，也可以直接验证：与反支配 Verma 模作
卷积给出等价
%
\[
  \Dmod(B \backslash G / B) \simeq \fg\mod_0^B,
\]
%
由此形式地推出 $\Gamma: \Dmod(\Fl_G) \rightarrow \fg\mod_0$ 的完全忠实性，参看
\cite{ahc}。
